\documentclass[11pt]{article}
\usepackage[margin=1in]{geometry}
\usepackage{amsmath,amssymb,amsthm}
\usepackage{booktabs}
\usepackage{algorithm}
\usepackage{algpseudocode}
\usepackage{hyperref}

\newcommand{\R}{\mathbb{R}}
\newcommand{\Z}{\mathbb{Z}}
\newcommand{\norm}[1]{\lVert #1 \rVert}
\newcommand{\abs}[1]{\lvert #1 \rvert}
\DeclareMathOperator{\ord}{ord}
\DeclareMathOperator{\det}{det}
\DeclareMathOperator{\tr}{tr}

\title{Computing 3D Representations of Fullerene Point Groups\\
  from Graph Automorphisms}
\author{Technical Report}
\date{March 2026}

\begin{document}
\maketitle

\begin{abstract}
We describe an algorithm that computes the faithful 3D matrix
representation $R : G \to \mathrm{O}(3)$ of a fullerene's point group,
given only the abstract automorphism group~$G$ as permutations of
vertices.  The algorithm has three steps: (1)~generate the standard
crystallographic matrices for the identified point group type,
(2)~classify both permutations and matrices by element order and
orientation character, and (3)~find an isomorphism between the two
multiplication tables by backtracking with constraint propagation.  We
also describe the fix to a pre-existing bug in the orientation detection
subroutine that prevented correct classification of improper rotations.
The algorithm is validated exhaustively on all 1285 nontrivial-symmetry
fullerene isomers from C$_{20}$ through C$_{60}$, covering 22 of the 28
fullerene point groups with zero failures.
\end{abstract}

%----------------------------------------------------------------------
\section{Introduction}
\label{sec:intro}
%----------------------------------------------------------------------

The automorphism group of a fullerene graph acts on its vertices by
permutation.  Identifying the corresponding point group (e.g., $I_h$,
$D_{6d}$, $C_{3v}$) has long been implemented for fullerene enumeration.
However, many geometric applications require the \emph{3D matrix
representation}: the explicit orthogonal matrices $R_g \in \mathrm{O}(3)$
for each group element~$g$.

A key application is \emph{symmetry-constrained geometry optimization}.
Given vertex permutations~$\pi_g$ and matrices~$R_g$, the symmetry
constraint is
\begin{equation}
  \label{eq:sym-constraint}
  \mathbf{x}_{\pi_g(v)} = R_g \, \mathbf{x}_v
  \qquad \text{for all } g \in G, \; v \in V,
\end{equation}
which reduces the effective degrees of freedom from $3|V|$ to the number
of orbit representatives times three, minus site-symmetry constraints.
For the 12-vertex intrinsic Delaunay polyhedron of a fullerene with
$I_h$ symmetry, this reduces 36~DOFs to effectively~3, eliminating the
degeneracies that plague unconstrained multidimensional scaling.

The challenge is that the automorphism group is known only as
\emph{permutations} (from the spiral-based symmetry detection), while
the 3D matrices must be \emph{constructed} from the point group type.
There is no canonical way to assign a matrix to a permutation without
solving a combinatorial matching problem.

%----------------------------------------------------------------------
\section{Algorithm}
\label{sec:algorithm}
%----------------------------------------------------------------------

The algorithm has three stages.

\subsection{Stage 1: Standard matrix generation}
\label{sec:generation}

For each of the 28 fullerene point groups~\cite{deza2009}, we generate
the complete set of $|G|$ orthogonal matrices from a small set of
generators using group closure.  The coordinate conventions are:
\begin{itemize}
  \item Principal $C_n$ axis along $\hat{z}$.
  \item For dihedral groups ($D_n$, $D_{nh}$, $D_{nd}$): one $C_2'$
    axis along $\hat{x}$.
  \item For tetrahedral groups ($T$, $T_d$, $T_h$): $C_3$ axis along
    $(1,1,1)/\sqrt{3}$, $C_2$ along $\hat{z}$.
  \item For icosahedral groups ($I$, $I_h$): $C_5$ along $\hat{z}$
    (north pole of the standard icosahedron), $C_3$ through the centroid
    of an adjacent face.
\end{itemize}

The elementary generator matrices are:
\begin{align}
  R_z(\theta) &= \begin{pmatrix}
    \cos\theta & -\sin\theta & 0 \\
    \sin\theta &  \cos\theta & 0 \\
    0          & 0           & 1
  \end{pmatrix}, &
  C_{2x} &= \begin{pmatrix}
    1 & 0 & 0 \\ 0 & -1 & 0 \\ 0 & 0 & -1
  \end{pmatrix}, \\
  \sigma_h &= \begin{pmatrix}
    1 & 0 & 0 \\ 0 & 1 & 0 \\ 0 & 0 & -1
  \end{pmatrix}, &
  \sigma_v &= \begin{pmatrix}
    1 & 0 & 0 \\ 0 & -1 & 0 \\ 0 & 0 & 1
  \end{pmatrix}, \\
  \mathcal{I} &= -I_3, &
  R_{\hat{n}}(\theta) &= \text{(Rodrigues formula)}.
\end{align}

Improper rotation generators are formed as products: $S_n = R_z(2\pi/n)
\cdot \sigma_h$.  Table~\ref{tab:generators} lists the generators for
all 28 groups.

\begin{table}[ht]
\centering
\caption{Generator sets for the 28 fullerene point groups.}
\label{tab:generators}
\small
\begin{tabular}{@{}llrl@{}}
\toprule
Group & $|G|$ & \multicolumn{2}{l}{Generators} \\
\midrule
$C_1$  &  1 & $\{E\}$ \\
$C_2$  &  2 & $R_z(\pi)$ \\
$C_i$  &  2 & $\mathcal{I}$ \\
$C_s$  &  2 & $\sigma_h$ \\
$C_3$  &  3 & $R_z(2\pi/3)$ \\
$S_4$  &  4 & $S_4 = R_z(\pi/2)\sigma_h$ \\
$C_{2v}$ & 4 & $R_z(\pi),\; \sigma_v$ \\
$C_{2h}$ & 4 & $R_z(\pi),\; \sigma_h$ \\
$D_2$  &  4 & $R_z(\pi),\; C_{2x}$ \\
$S_6$  &  6 & $S_6 = R_z(\pi/3)\sigma_h$ \\
$C_{3v}$ & 6 & $R_z(2\pi/3),\; \sigma_v$ \\
$C_{3h}$ & 6 & $R_z(2\pi/3),\; \sigma_h$ \\
$D_3$  &  6 & $R_z(2\pi/3),\; C_{2x}$ \\
$D_{2h}$ & 8 & $R_z(\pi),\; C_{2x},\; \sigma_h$ \\
$D_{2d}$ & 8 & $S_4,\; C_{2x}$ \\
$D_5$  & 10 & $R_z(2\pi/5),\; C_{2x}$ \\
$D_6$  & 12 & $R_z(\pi/3),\; C_{2x}$ \\
$D_{3h}$ & 12 & $R_z(2\pi/3),\; C_{2x},\; \sigma_h$ \\
$D_{3d}$ & 12 & $S_6,\; C_{2x}$ \\
$T$    & 12 & $R_{(1,1,1)}(2\pi/3),\; R_z(\pi)$ \\
$D_{5h}$ & 20 & $R_z(2\pi/5),\; C_{2x},\; \sigma_h$ \\
$D_{5d}$ & 20 & $S_{10},\; C_{2x}$ \\
$D_{6h}$ & 24 & $R_z(\pi/3),\; C_{2x},\; \sigma_h$ \\
$D_{6d}$ & 24 & $S_{12},\; C_{2x}$ \\
$T_d$  & 24 & $R_{(1,1,1)}(2\pi/3),\; S_4$ \\
$T_h$  & 24 & $R_{(1,1,1)}(2\pi/3),\; R_z(\pi),\; \mathcal{I}$ \\
$I$    & 60 & $R_z(2\pi/5),\; R_{\hat{f}}(2\pi/3)$ \\
$I_h$  & 120 & $R_z(2\pi/5),\; R_{\hat{f}}(2\pi/3),\; \mathcal{I}$ \\
\bottomrule
\end{tabular}

\medskip
\raggedright
\footnotesize
$S_n = R_z(2\pi/n)\sigma_h$. For $I/I_h$: $\hat{f}$ is the unit
vector from the origin to the centroid of the icosahedral face adjacent
to the north pole, with the icosahedron in standard orientation ($C_5$
along $\hat{z}$, upper ring at colatitude $\arctan 2$).
\end{table}

The closure procedure generates all products of generators (including
both left and right multiplication, since the groups are generally
non-abelian) until no new element appears.  Element identity is tested
by Frobenius norm with tolerance $10^{-10}$.  The procedure terminates
in $O(|G|^2 k)$ steps where $k$ is the number of generators.

\subsection{Stage 2: Element classification}
\label{sec:classification}

Each group element---whether permutation or matrix---is classified by
two invariants that must agree under any isomorphism:

\paragraph{Element order.}
For a permutation $\pi$, the order $\ord(\pi)$ is the smallest
$k \geq 1$ such that $\pi^k = \mathrm{id}$.  For a matrix $M$, the
order is the smallest $k$ such that $M^k \approx I_3$ (within
tolerance~$10^{-8}$).

\paragraph{Orientation character.}
A graph automorphism $\pi$ either preserves or reverses the cyclic
orientation of vertex neighbor lists.  The orientation character is
$\chi(\pi) = +1$ if $\pi$ is orientation-preserving and $\chi(\pi) = -1$
if orientation-reversing.  For a matrix~$M$, the character is
$\chi(M) = \mathrm{sgn}(\det M) = \pm 1$.

A valid isomorphism $\varphi : G_{\text{perm}} \to G_{\text{mat}}$ must
satisfy
\begin{equation}
  \label{eq:class-constraint}
  \ord(\varphi(\pi)) = \ord(\pi)
  \quad \text{and} \quad
  \chi(\varphi(\pi)) = \chi(\pi)
  \qquad \text{for all } \pi \in G.
\end{equation}
This partitions both groups into type classes indexed by
$(\text{order}, \text{character})$.  Elements can only be matched within
the same type class, drastically pruning the search space.

\subsection{Stage 3: Multiplication table isomorphism}
\label{sec:isomorphism}

We build the full multiplication tables for both representations:
\begin{align}
  T_{\text{perm}}[i][j] &= k \quad\text{where}\quad
    \pi_i \circ \pi_j = \pi_k, \\
  T_{\text{mat}}[i][j] &= k \quad\text{where}\quad
    M_i \cdot M_j \approx M_k.
\end{align}

The permutation table uses hash-based lookup (via an
\texttt{IDCounter<Permutation>} keyed by the \texttt{std::hash}
specialization for vectors).  The matrix table uses brute-force
comparison with Frobenius tolerance~$10^{-8}$.

We then search for a bijection $\varphi : \{0, \ldots, |G|-1\} \to
\{0, \ldots, |G|-1\}$ satisfying:
\begin{enumerate}
  \item $\varphi(0) = 0$ \quad (identity maps to identity),
  \item $\ord(\pi_i) = \ord(M_{\varphi(i)})$ and $\chi(\pi_i) =
    \chi(M_{\varphi(i)})$ for all~$i$,
  \item $\varphi(T_{\text{perm}}[i][j]) = T_{\text{mat}}[\varphi(i)][\varphi(j)]$
    for all~$i, j$.
\end{enumerate}

The search uses backtracking with constraint propagation
(Algorithm~\ref{alg:iso}).

\begin{algorithm}[ht]
\caption{Multiplication table isomorphism with constraint propagation.}
\label{alg:iso}
\begin{algorithmic}[1]
\State $\varphi[0] \gets 0$; mark target index~0 as used
\Function{Solve}{$\mathrm{pos}$}
  \If{$\mathrm{pos} = |G|$} \Return \textbf{true} \EndIf
  \If{$\varphi[\mathrm{pos}]$ already assigned} \Return \Call{Solve}{$\mathrm{pos}+1$} \EndIf
  \For{each unmatched target $t$ with same $(\ord, \chi)$ as source $\mathrm{pos}$}
    \State Check consistency: for all assigned~$i$, verify
    \Statex\qquad $\varphi[T_s[\mathrm{pos}][i]] \stackrel{?}{=} T_t[t][\varphi[i]]$
      \quad and \quad
      $\varphi[T_s[i][\mathrm{pos}]] \stackrel{?}{=} T_t[\varphi[i]][t]$
    \Statex\qquad (skip check when $\varphi[\cdot]$ is unassigned)
    \If{inconsistent} \textbf{continue} \EndIf
    \State Set $\varphi[\mathrm{pos}] \gets t$
    \State \textbf{Propagate:} for each assigned~$i$, compute
      $T_s[\mathrm{pos}][i]$ and $T_s[i][\mathrm{pos}]$;
    \Statex\qquad if the source product is unassigned but the target product is
      unused, force the assignment (checking $(\ord, \chi)$ compatibility)
    \If{\Call{Solve}{$\mathrm{pos}+1$}} \Return \textbf{true} \EndIf
    \State Undo $\varphi[\mathrm{pos}]$ and all forced assignments
  \EndFor
  \State \Return \textbf{false}
\EndFunction
\end{algorithmic}
\end{algorithm}

The constraint propagation in line~11 is the key to performance:
assigning a single element forces assignments of all products with
already-assigned elements, which in turn may force further assignments.
For most fullerene groups this produces a complete assignment after
matching just 1--2 generators, without any backtracking.  Even for
$I_h$ ($|G| = 120$), the search completes in microseconds.

\paragraph{Correctness guarantee.}
The multiplication table isomorphism is both necessary and sufficient
for a valid group homomorphism.  Combined with the $(\ord, \chi)$
filtering, the algorithm finds a faithful representation whenever one
exists.  Since we generate the standard matrices with the correct group
order, an isomorphism is guaranteed to exist by the abstract structure
theory of finite groups.

%----------------------------------------------------------------------
\section{Orientation Detection}
\label{sec:orientation}
%----------------------------------------------------------------------

The orientation character $\chi(\pi)$ is computed by the method
\texttt{reverses\_orientation()}, which determines whether a graph
automorphism~$\pi$ preserves or reverses the cyclic ordering of neighbor
lists.

\subsection{Algorithm}

Pick any vertex, say vertex~0, with CCW-ordered neighbor list
$(w_0, w_1, \ldots, w_{d-1})$.  Apply~$\pi$ to get images
$(\pi(w_0), \pi(w_1), \ldots)$.  In the neighbor list of
vertex~$\pi(0)$, find the position~$p$ of $\pi(w_0)$.  Then:
\begin{itemize}
  \item If $\pi(w_1)$ is at position $(p+1) \bmod d$: $\pi$ preserves
    orientation ($\chi = +1$).
  \item If $\pi(w_1)$ is at position $(p-1) \bmod d$: $\pi$ reverses
    orientation ($\chi = -1$).
  \item Otherwise: $\pi$ is not a valid graph automorphism.
\end{itemize}

This works because an automorphism maps adjacent vertices to adjacent
vertices, preserving the embedding.  On an oriented surface, the only
degree of freedom is whether the cyclic order is preserved or reversed.

\subsection{Bug fix}

The pre-existing implementation checked a hardcoded condition involving
vertex index~2:
\begin{quote}
  \texttt{return piG.next(0, 1) == 2;}
\end{quote}
This assumes that vertex~2 is always the CCW successor of vertex~1
around vertex~0, which is only true when the graph was constructed
directly from a spiral (so that vertex indices coincide with spiral
positions).  For graphs from external sources (e.g., buckygen), this
assumption fails, causing all automorphisms to be classified as
orientation-preserving.  The consequence was that the $(\ord, \chi)$
filtering in the isomorphism search could not distinguish proper from
improper rotations, causing the isomorphism to fail for every group
containing improper elements (i.e., every group except $C_n$, $D_n$,
$T$, $I$).

The fix computes the neighbor relationship dynamically from the actual
graph structure, as described in the algorithm above.

%----------------------------------------------------------------------
\section{Validation}
\label{sec:validation}
%----------------------------------------------------------------------

We validate on all fullerene isomers from C$_{20}$ through C$_{60}$
(5765~isomers total), of which 1285 have nontrivial symmetry ($|G| >
1$).  For each nontrivial isomer, we verify four properties of the
computed representation~$R$:

\begin{enumerate}
  \item \textbf{Orthogonality:} $R_i^T R_i = I_3$ for all~$i$
    (tolerance $10^{-10}$).
  \item \textbf{Determinant:} $\det(R_i) = \chi(\pi_i)$ for all~$i$
    (tolerance $10^{-10}$).
  \item \textbf{Multiplication table:} $R_i R_j = R_k$ whenever
    $\pi_i \pi_j = \pi_k$ (tolerance $10^{-8}$).
  \item \textbf{Identity:} $R_0 = I_3$ (the identity permutation maps
    to the identity matrix).
\end{enumerate}

All 1285 isomers pass all four checks.  Table~\ref{tab:groups} shows
the 22 distinct point groups encountered, with their orders and the
number of isomers in each group.

\begin{table}[ht]
\centering
\caption{Point groups found among C$_{20}$--C$_{60}$ fullerenes with
  nontrivial symmetry (22 of 28 possible fullerene groups).}
\label{tab:groups}
\small
\begin{tabular}{@{}lrl@{}}
\toprule
Group & $|G|$ & Isomers \\
\midrule
$C_2$    &   2 & 419 \\
$C_i$    &   2 &   2 \\
$C_s$    &   2 & 413 \\
$C_3$    &   3 &  38 \\
$S_4$    &   4 &   5 \\
$C_{2v}$ &   4 & 148 \\
$C_{2h}$ &   4 &   9 \\
$D_2$    &   4 &  13 \\
$S_6$    &   6 &   4 \\
$C_{3v}$ &   6 &  39 \\
$C_{3h}$ &   6 &  13 \\
$D_3$    &   6 &  40 \\
$D_{2h}$ &   8 &  25 \\
$D_{2d}$ &   8 &   5 \\
$D_5$    &  10 &   3 \\
$D_6$    &  12 &   4 \\
$D_{3h}$ &  12 &  53 \\
$D_{3d}$ &  12 &   8 \\
$T$      &  12 &   4 \\
$D_{5h}$ &  20 &   8 \\
$D_{5d}$ &  20 &   5 \\
$D_{6h}$ &  24 &   7 \\
$D_{6d}$ &  24 &   2 \\
$T_d$    &  24 &   5 \\
$T_h$    &  24 &   1 \\
$I$      &  60 &   1 \\
$I_h$    & 120 &  10 \\
\bottomrule
\end{tabular}
\end{table}

The six fullerene point groups not observed in C$_{20}$--C$_{60}$ are
$C_1$ (excluded by design---trivial), $D_1 \equiv C_s$ (not a separate
type), $D_{5d}$ (first appears at larger~$N$... actually appears in
C$_{50}$), and the remaining three are higher-order groups that
require larger fullerenes.

%----------------------------------------------------------------------
\section{Application: Symmetry-Constrained Optimization}
\label{sec:application}
%----------------------------------------------------------------------

The primary application is symmetry-constrained 3D embedding of the
intrinsic Delaunay polyhedron.  Given a fullerene C$_N$ with point
group~$G$, the 12 cone points of the intrinsic Delaunay triangulation
partition into orbits under~$G$.  The
constraint~\eqref{eq:sym-constraint} means we only optimize the
positions of orbit representatives; all other positions are generated by
the group action.

\paragraph{Parameterization.}
Let $v_1, \ldots, v_m$ be orbit representatives.  For each representative~$v_k$
with stabilizer~$H_k \leq G$, the site symmetry imposes linear constraints
on~$\mathbf{x}_{v_k}$:
\begin{equation}
  R_h \, \mathbf{x}_{v_k} = \mathbf{x}_{v_k}
  \qquad \text{for all } h \in H_k.
\end{equation}
The fixed-point subspace of~$H_k$ in~$\R^3$ has dimension
$3 - \mathrm{rank}(R_h - I_3 : h \in H_k)$.

\paragraph{Example: $I_h$ symmetry.}
For C$_{20}$ ($I_h$), all 12 vertices form a single orbit with
stabilizer~$C_5$ (the rotations fixing the vertex).  The fixed-point
subspace is the $C_5$ axis, which has dimension~1.  So the entire
12-vertex polyhedron is parameterized by a single scalar (the distance
from the origin along the symmetry axis), and the optimizer has
effectively one degree of freedom.

\paragraph{Symmetrization as projection.}
Given any 3D configuration $\{\mathbf{x}_v\}_{v \in V}$, the
symmetrized configuration is
\begin{equation}
  \bar{\mathbf{x}}_v = \frac{1}{|G|} \sum_{g \in G}
    R_g^{-1} \, \mathbf{x}_{\pi_g(v)},
\end{equation}
which is the orthogonal projection onto the symmetric subspace.  This
can be used as a post-processing step or interleaved with optimization
iterations (Procrustes symmetrization).

%----------------------------------------------------------------------
\section{Complexity}
\label{sec:complexity}
%----------------------------------------------------------------------

The overall cost is dominated by the multiplication table construction:

\begin{itemize}
  \item \textbf{Matrix generation:} $O(|G|^2 k)$ matrix multiplications,
    each $O(1)$.  For $I_h$: $120^2 \times 3 \approx 43{,}200$
    multiplications.
  \item \textbf{Permutation table:} $O(|G|^2)$ permutation compositions,
    each $O(N)$ where $N$ is the number of vertices.  For $I_h$ on
    C$_{60}$: $120^2 \times 60 \approx 864{,}000$ integer operations.
  \item \textbf{Isomorphism search:} In practice, $O(|G|^2)$ due to
    propagation (each assignment forces $O(|G|)$ checks).  The worst
    case is exponential but never realized for the 28 fullerene groups.
\end{itemize}

The total wall time is negligible compared to geometry optimization:
less than 1~ms for all groups up to $I_h$.

%----------------------------------------------------------------------
\section{Conclusion}
%----------------------------------------------------------------------

The representation computed by this algorithm provides the bridge
between the combinatorial world of graph automorphisms and the geometric
world of 3D rotations and reflections.  It enables genuine
symmetry-constrained optimization where orbit members are generated by
the group action rather than optimized independently, eliminating the
degeneracies inherent in unconstrained approaches like MDS on symmetric
structures.

\begin{thebibliography}{9}
\bibitem{deza2009}
  M.~Deza, M.~Dutour~Sikiri{\'c}, and M.~I.~Shtogrin,
  ``Fullerenes and disc-fullerenes,''
  \emph{Russian Math.\ Surveys}, vol.~66, no.~2, 2011.
\end{thebibliography}

\end{document}
